\chapter{绪论}

\section{研究背景}

计算成像系统正在从“先光学设计、后算法补偿”的串行范式逐步转向光学与算法协同的联合设计范式。
在移动成像、工业检测和智能感知等应用场景中，系统同时面临体积受限、视场提升和成像质量稳定的多目标约束。
传统纯折射系统在小型化与宽视场条件下容易引入像差累积，而单纯依赖后处理会放大噪声并损伤细节可解释性。

\section{研究问题与目标}

本论文聚焦于三个核心问题：
第一，如何在统一可微框架中同时刻画折射传播与衍射传播；
第二，如何在高维耦合优化中保证训练稳定性和可复现性；
第三，如何建立满足学术审查要求的定量评估与证据闭环。

对应目标为：构建可执行的端到端折衍混合设计流程，形成可复核的实验体系，并给出工程可落地的参数配置建议。
